(1) 粒子从静止点沿重力场中光滑曲线下滑，纵向坐标 \(y\) 向下为正，证明下降时间泛函 \(T=\int \sqrt{1+y'^2}\,dx/\sqrt{2gy}\)。
(2) 由Beltrami恒等式求最速降线的首次积分，并用参数 \(\theta\) 给出摆线方程。
(3) 若粒子沿竖直平面倒摆线槽运动，最低点附近向左右摆动，证明由弧长坐标描述的势能是严格二次型。
(4) 推导从槽壁任意初始高度静止释放到最低点的时间，说明何谓等时性。
